%%%%%%%%%%%%%%%%%%%%%%%%%%%%%%%%%%%%%%%%%%%%%%%%%%%%%%%%%%%
% Slide: Course Overview
%%%%%%%%%%%%%%%%%%%%%%%%%%%%%%%%%%%%%%%%%%%%%%%%%%%%%%%%%%%
\begin{frame}[fragile]\frametitle{Course Overview}

      \begin{itemize}
        \item Python Fundamentals
        \item Object-Oriented Programming
        \item File Handling \& I/O Operations
        \item Advanced Python Features
        \item Asynchronous Programming
        \item Real-world Projects
      \end{itemize}

\end{frame}

%%%%%%%%%%%%%%%%%%%%%%%%%%%%%%%%%%%%%%%%%%%%%%%%%%%%%%%%%%%
% Slide: Context Managers
%%%%%%%%%%%%%%%%%%%%%%%%%%%%%%%%%%%%%%%%%%%%%%%%%%%%%%%%%%%
\begin{frame}[fragile]\frametitle{Context Managers}
\begin{columns}
    \begin{column}[T]{0.5\linewidth}
      \begin{itemize}
        \item with statement for resource management
        \item \_\_enter\_\_ and \_\_exit\_\_ methods
        \item contextlib module
        \item @contextmanager decorator
        \item Automatic cleanup of resources
        \item Custom context managers
      \end{itemize}
    \end{column}
    \begin{column}[T]{0.5\linewidth}
      \begin{lstlisting}[language=python, basicstyle=\tiny]
from contextlib import contextmanager

# Custom context manager class
class FileManager:
    def __init__(self, filename, mode):
        self.filename = filename
        self.mode = mode

    def __enter__(self):
        self.file = open(self.filename, self.mode)
        return self.file

    def __exit__(self, exc_type, exc_val, exc_tb):
        self.file.close()

# Using decorator
@contextmanager
def timer_context():
    start = time.time()
    yield
    print(f"Elapsed: {time.time()-start:.2f}s")

with timer_context():
    # Code to time
    time.sleep(1)
      \end{lstlisting}
    \end{column}
  \end{columns}
\end{frame}

%%%%%%%%%%%%%%%%%%%%%%%%%%%%%%%%%%%%%%%%%%%%%%%%%%%%%%%%%%%
% Slide: Threading & Multiprocessing
%%%%%%%%%%%%%%%%%%%%%%%%%%%%%%%%%%%%%%%%%%%%%%%%%%%%%%%%%%%
\begin{frame}[fragile]\frametitle{Concurrent Operations - Threading}
\begin{columns}
    \begin{column}[T]{0.5\linewidth}
      \begin{itemize}
        \item threading module for concurrency
        \item Thread creation and management
        \item Global Interpreter Lock (GIL)
        \item Thread-safe operations
        \item concurrent.futures.ThreadPoolExecutor
        \item Use for I/O-bound tasks
      \end{itemize}
    \end{column}
    \begin{column}[T]{0.5\linewidth}
      \begin{lstlisting}[language=python, basicstyle=\tiny]
import threading
from concurrent.futures import ThreadPoolExecutor
import time

# Basic threading
def task(name):
    print(f"Task {name} starting")
    time.sleep(1)
    print(f"Task {name} complete")

threads = []
for i in range(3):
    t = threading.Thread(target=task, args=(i,))
    threads.append(t)
    t.start()

for t in threads:
    t.join()

# ThreadPoolExecutor
with ThreadPoolExecutor(max_workers=3) as executor:
    futures = [executor.submit(task, i) for i in range(3)]
      \end{lstlisting}
    \end{column}
  \end{columns}
\end{frame}

%%%%%%%%%%%%%%%%%%%%%%%%%%%%%%%%%%%%%%%%%%%%%%%%%%%%%%%%%%%
% Slide: Async/Await Basics
%%%%%%%%%%%%%%%%%%%%%%%%%%%%%%%%%%%%%%%%%%%%%%%%%%%%%%%%%%%
\begin{frame}[fragile]\frametitle{Async/Await Basics}
\begin{columns}
    \begin{column}[T]{0.5\linewidth}
      \begin{itemize}
        \item Asynchronous programming concepts
        \item async def for coroutines
        \item await keyword for awaiting
        \item Event loop with asyncio
        \item Concurrent execution of coroutines
        \item Non-blocking I/O operations
      \end{itemize}
    \end{column}
    \begin{column}[T]{0.5\linewidth}
      \begin{lstlisting}[language=python, basicstyle=\tiny]
import asyncio

# Basic async function
async def fetch_data(id):
    print(f"Fetching {id}...")
    await asyncio.sleep(1)  # Simulates I/O
    print(f"Fetched {id}")
    return f"Data {id}"

# Running coroutines
async def main():
    # Sequential
    result1 = await fetch_data(1)

    # Concurrent
    results = await asyncio.gather(
        fetch_data(2),
        fetch_data(3),
        fetch_data(4)
    )
    print(results)

# Run event loop
asyncio.run(main())
      \end{lstlisting}
    \end{column}
  \end{columns}
\end{frame}

%%%%%%%%%%%%%%%%%%%%%%%%%%%%%%%%%%%%%%%%%%%%%%%%%%%%%%%%%%%
% Slide: Project - Library Management System
%%%%%%%%%%%%%%%%%%%%%%%%%%%%%%%%%%%%%%%%%%%%%%%%%%%%%%%%%%%
\begin{frame}[fragile]\frametitle{Project: Library Management System}
\begin{lstlisting}[language=python, basicstyle=\tiny]
class Book:
    def __init__(self, title, author, isbn):
        self.title = title
        self.author = author
        self.isbn = isbn
        self.is_available = True

class Member:
    def __init__(self, name, member_id):
        self.name = name
        self.member_id = member_id
        self.borrowed_books = []

class Library:
    def __init__(self):
        self.books = []
        self.members = []

    def add_book(self, book): # TODO: Implement adding book to library
        pass

    def lend_book(self, member, book): # TODO: Check availability, update records
        pass

    def return_book(self, member, book): # TODO: Process book return
        pass
\end{lstlisting}
\end{frame}

%%%%%%%%%%%%%%%%%%%%%%%%%%%%%%%%%%%%%%%%%%%%%%%%%%%%%%%%%%%
% Slide: Project - Shopping Cart System
%%%%%%%%%%%%%%%%%%%%%%%%%%%%%%%%%%%%%%%%%%%%%%%%%%%%%%%%%%%
\begin{frame}[fragile]\frametitle{Project: Shopping Cart System}
\begin{lstlisting}[language=python, basicstyle=\tiny]
class Product:
    def __init__(self, name, price, category):
        self.name = name
        self.price = price
        self.category = category

class ShoppingCart:
    def __init__(self):
        self.items = {}  # {product: quantity}

    def add_item(self, product, quantity=1): # TODO: Add product to cart
        pass

    def remove_item(self, product): # TODO: Remove product from cart
        pass

    def calculate_total(self, discount_code=None):
        # TODO: Calculate total with optional discount
        # Apply percentage discount if code valid
        pass

    def checkout(self): # TODO: Process payment and clear cart
        pass
\end{lstlisting}
\end{frame}

%%%%%%%%%%%%%%%%%%%%%%%%%%%%%%%%%%%%%%%%%%%%%%%%%%%%%%%%%%%
% Slide: Project - CSV Data Processor
%%%%%%%%%%%%%%%%%%%%%%%%%%%%%%%%%%%%%%%%%%%%%%%%%%%%%%%%%%%
\begin{frame}[fragile]\frametitle{Project: CSV Data Processor}
\begin{lstlisting}[language=python, basicstyle=\tiny]
import csv
import logging

def log_operation(func):
    """Decorator to log operations"""
    def wrapper(*args, **kwargs):
        logging.info(f"Starting {func.__name__}")
        try:
            result = func(*args, **kwargs)
            logging.info(f"Completed {func.__name__}")
            return result
        except Exception as e:
            logging.error(f"Error in {func.__name__}: {e}")
            raise
    return wrapper

class CSVProcessor:
    def __init__(self, input_file):
        self.input_file = input_file
    @log_operation
    def read_data(self): # TODO: Read CSV with error handling
        pass
    @log_operation
    def filter_data(self, condition): # TODO: Filter rows based on condition
        pass
    @log_operation
    def transform_data(self, transformation): # TODO: Apply transformation to columns
        pass
    @log_operation
    def export_results(self, output_file): # TODO: Write processed data to new CSV
        pass
\end{lstlisting}
\end{frame}

%%%%%%%%%%%%%%%%%%%%%%%%%%%%%%%%%%%%%%%%%%%%%%%%%%%%%%%%%%%
% Slide: Project - Config File Manager
%%%%%%%%%%%%%%%%%%%%%%%%%%%%%%%%%%%%%%%%%%%%%%%%%%%%%%%%%%%
\begin{frame}[fragile]\frametitle{Project: Config File Manager}
\begin{lstlisting}[language=python, basicstyle=\tiny]
import json
import yaml
from contextlib import contextmanager
from datetime import datetime

class ConfigManager:
    def __init__(self, config_file):
        self.config_file = config_file
        self.config = {}

    @contextmanager
    def config_context(self): # TODO: Backup current config
        """Context manager for safe config operations"""
        try:
            yield self
        except Exception as e: # TODO: Restore from backup
            raise
        finally: # TODO: Cleanup
            pass

    def load_config(self): # TODO: Load JSON/YAML config
        pass
    def validate_config(self, schema): # TODO: Validate against schema
        pass
    def save_config(self): # TODO: Save config with backup
        pass
    def get(self, key, default=None): # TODO: Get config value with dot notation
        pass
\end{lstlisting}
\end{frame}

%%%%%%%%%%%%%%%%%%%%%%%%%%%%%%%%%%%%%%%%%%%%%%%%%%%%%%%%%%%
% Slide: Project - Async Web Scraper
%%%%%%%%%%%%%%%%%%%%%%%%%%%%%%%%%%%%%%%%%%%%%%%%%%%%%%%%%%%
\begin{frame}[fragile]\frametitle{Project: Async Web Scraper}
\begin{lstlisting}[language=python, basicstyle=\tiny]
import asyncio
import aiohttp
from bs4 import BeautifulSoup

class AsyncWebScraper:
    def __init__(self, urls):
        self.urls = urls
        self.results = []

    async def fetch_url(self, session, url):
        """Fetch single URL asynchronously"""
        try: # TODO: Get HTML content, Parse with BeautifulSoup, Extract required data
            async with session.get(url) as response:
                pass
        except Exception as e:
            print(f"Error fetching {url}: {e}")
            return None

    async def scrape_all(self): # TODO: Gather all results concurrently
        async with aiohttp.ClientSession() as session:
            tasks = [self.fetch_url(session, url) for url in self.urls]
            pass

    async def save_results(self, filename): # TODO: Write results to file
        pass

# asyncio.run(scraper.scrape_all())
\end{lstlisting}
\end{frame}

%%%%%%%%%%%%%%%%%%%%%%%%%%%%%%%%%%%%%%%%%%%%%%%%%%%%%%%%%%%
% Slide: Project - Concurrent File Processor
%%%%%%%%%%%%%%%%%%%%%%%%%%%%%%%%%%%%%%%%%%%%%%%%%%%%%%%%%%%
\begin{frame}[fragile]\frametitle{Project: Concurrent File Processor}
\begin{lstlisting}[language=python, basicstyle=\tiny]
import asyncio
import aiofiles
from pathlib import Path

class ConcurrentFileProcessor:
    def __init__(self, file_paths):
        self.file_paths = file_paths
        self.progress = {}

    async def process_file(self, file_path):
        try: # TODO: Read file in chunks, Process, Update progress
            async with aiofiles.open(file_path, 'r') as f:
                pass
        except Exception as e:
            print(f"Error processing {file_path}: {e}")

    async def track_progress(self):
        while True: # TODO: Display progress for all files
            await asyncio.sleep(0.5) # TODO: Break when all complete
            pass

    async def process_all(self): # TODO: Run all tasks concurrently
        tasks = [self.process_file(fp) for fp in self.file_paths]
        progress_task = asyncio.create_task(self.track_progress())
        pass

# processor = ConcurrentFileProcessor(file_list)
# asyncio.run(processor.process_all())
\end{lstlisting}
\end{frame}
